% block5_trinomial.tex
%
% Phase 3, Block 5: Trinomial tree of Kamrad-Ritchken.
%
% Self-contained writeup. Compiles standalone with pdflatex.

\documentclass[11pt,a4paper]{article}

\usepackage[T1]{fontenc}
\usepackage[utf8]{inputenc}
\usepackage{amsmath, amssymb, amsthm}
\usepackage{geometry}
\usepackage{hyperref}
\usepackage{enumitem}

\geometry{margin=1in}

\newtheorem{proposition}{Proposition}[section]
\newtheorem{theorem}[proposition]{Theorem}
\newtheorem{remark}[proposition]{Remark}

\title{Phase 3, Block 5:\\The Kamrad--Ritchken trinomial tree\\and its equivalence to FTCS}
\author{Quant Finance Portfolio}
\date{}

\begin{document}
\maketitle

\section{Setting and motivation}

Phase 2 Block 6 already implemented a CRR binomial tree for American
options. Phase 3 has built four PDE-based pricers (FTCS, BTCS, CN,
and CN-PSOR for American). What does the trinomial tree add?

Not pricing capability: the binomial CRR already prices Americans;
CN-PSOR already prices them more accurately. The contribution of
Block 5 is conceptual --- it makes \emph{explicit the structural
equivalence between lattice methods and finite-difference
methods}. We have alluded to this equivalence in earlier blocks
(``FTCS and the trinomial are the same calculation written in two
languages''); Block 5 formalises it.

Two additional benefits:
\begin{itemize}[leftmargin=2em]
\item Block 6 will benchmark all the project's pricers on the
American put. The trinomial gives one more independent point of
comparison.
\item The trinomial of Kamrad and Ritchken~\cite{KR1991} introduces
a stretching parameter $\lambda$ that does not exist in the
binomial. This adds a tuning knob to the family of lattice methods
which will reappear in Block~6's comparison.
\end{itemize}

% =====================================================================
\section{Tree structure}
% =====================================================================

We work in log-coordinates $x = \ln(S/S_0)$, so the GBM
$dS = rS\,dt + \sigma S\,dW$ becomes
\[
dX = \nu\,dt + \sigma\,dW, \qquad \nu = r - \tfrac{1}{2}\sigma^2.
\]
The trinomial tree discretises this process with time step $\Delta t = T/N$
and space step $\Delta x$ (to be chosen). At each node, the next
state branches three ways:
\[
x \;\longrightarrow\;
\begin{cases}
x + \Delta x & \text{with probability } p_u \\
x           & \text{with probability } p_m \\
x - \Delta x & \text{with probability } p_d
\end{cases}
\]
with $p_u + p_m + p_d = 1$. The key structural features:
\begin{itemize}[leftmargin=2em]
\item \textbf{Recombining}: an up-then-down move arrives at the same
state as a stay-then-stay move. After $n$ time steps, only $2n+1$
distinct states are reachable (not $3^n$). This is what makes the
backward induction tractable: $O(N^2)$ total work, not $O(3^N)$.
\item \textbf{Symmetric in space}: the same $\Delta x$ for up and
down. This simplifies the moment-matching algebra and makes the tree
agree with the centred-difference FD scheme of Block~1, as we will
see.
\end{itemize}

\paragraph{The Kamrad--Ritchken step size.} The KR specification picks
\begin{equation}
\Delta x = \sigma\sqrt{\lambda\,\Delta t}, \qquad \lambda \ge 1.
\label{eq:kr-step}
\end{equation}
The parameter $\lambda$ is the \emph{stretching parameter}.
Three values are commonly used:
\begin{center}
\begin{tabular}{ccc}
\hline
$\lambda$ & Tree type & $p_m$ \\
\hline
$1$ & Binomial-equivalent ($p_m = 0$) & $0$ \\
$2$ & Conservative trinomial          & $1/2$ \\
$3$ & Standard KR (``third-order'')   & $2/3$ \\
\hline
\end{tabular}
\end{center}
We will see in \S\ref{sec:equivalence} that $\lambda$ is structurally
the inverse of the FTCS Fourier number: $\lambda = 1/(2\alpha)$.
The constraint $\lambda \ge 1$ is exactly the FTCS CFL condition
$\alpha \le 1/2$ written as a tree constraint.

% =====================================================================
\section{Probabilities by moment matching}
% =====================================================================

The probabilities $(p_u, p_m, p_d)$ are determined by matching
the first two moments of the discrete one-step distribution to those
of the continuous GBM:
\begin{align}
\mathbb{E}[\Delta X] \;&=\; p_u\,\Delta x - p_d\,\Delta x \;=\; \nu\,\Delta t,
\label{eq:mom1}\\
\mathbb{E}[(\Delta X)^2] \;&=\; (p_u + p_d)\,\Delta x^2 \;=\; \sigma^2\,\Delta t + O(\Delta t^2).
\label{eq:mom2}
\end{align}
Combined with the normalisation $p_u + p_m + p_d = 1$, this is a
linear $3 \times 3$ system. Substituting \eqref{eq:kr-step} and solving:
\begin{equation}
\boxed{\;\;
p_u = \frac{1}{2\lambda} + \frac{\nu\sqrt{\Delta t}}{2\sigma\sqrt{\lambda}},
\quad
p_m = 1 - \frac{1}{\lambda},
\quad
p_d = \frac{1}{2\lambda} - \frac{\nu\sqrt{\Delta t}}{2\sigma\sqrt{\lambda}}.
\;\;}
\label{eq:probs}
\end{equation}

\paragraph{Reading.} The first term $1/(2\lambda)$ in $p_u$ and $p_d$
encodes the variance of the step. The second term
$\pm\nu\sqrt{\Delta t}/(2\sigma\sqrt{\lambda})$ is the asymmetry that
accommodates the drift. The middle probability $p_m$ is determined
purely by $\lambda$.

\paragraph{Positivity constraints.}
\begin{itemize}[leftmargin=2em]
\item $p_m \ge 0 \iff \lambda \ge 1$.
\item $p_d \ge 0 \iff \nu\sqrt{\Delta t} \le \sigma\sqrt{\lambda} \iff \dfrac{\nu^2\,\Delta t}{\sigma^2} \le \lambda$. For project parameters ($|\nu| \sim 0.03$, $\sigma = 0.20$, $\Delta t \le 0.02$) the LHS is $\le 1.5\times 10^{-4}$, far below any $\lambda \ge 1$.
\item $p_u, p_m, p_d \le 1$ is then automatic.
\end{itemize}
For the project's parameters, all three probabilities lie comfortably in
$(0, 1)$.

% =====================================================================
\section{Backward induction}
% =====================================================================

With probabilities fixed, the algorithm is standard.

\paragraph{Initialisation at maturity ($n = N$).} The states at time
$T$ are $S_j = S_0\,e^{j\,\Delta x}$ for $j \in \{-N, -N+1, \ldots, N\}$.
For each, set $V_j^N = \phi(S_j)$ where $\phi$ is the payoff function
($(K - S)^+$ for a put, $(S - K)^+$ for a call).

\paragraph{Backward step ($n = N-1, \ldots, 0$).} For each
$j \in \{-n, \ldots, n\}$,
\begin{equation}
V_j^n = e^{-r\Delta t}\bigl[\,p_u\,V_{j+1}^{n+1} + p_m\,V_j^{n+1} + p_d\,V_{j-1}^{n+1}\,\bigr]
\label{eq:european-step}
\end{equation}
for the European option, or
\begin{equation}
V_j^n = \max\!\Bigl\{\,\phi(S_j),\;\;
e^{-r\Delta t}\bigl[\,p_u\,V_{j+1}^{n+1} + p_m\,V_j^{n+1} + p_d\,V_{j-1}^{n+1}\,\bigr]\,\Bigr\}
\label{eq:american-step}
\end{equation}
for the American option. The price is $V_0^0$.

\paragraph{Cost.} At step $n$ there are $2n+1$ active nodes. Total
work is $\sum_{n=0}^{N-1}(2n+1) = N^2$ multiply-add operations, i.e.\
$O(N^2)$ time. With in-place overwrite of a single length-$(2N+1)$
array, memory is $O(N)$. Identical to CRR.

% =====================================================================
\section{The equivalence with FTCS}
\label{sec:equivalence}
% =====================================================================

This is the structural insight that motivates including the trinomial
in Phase 3. The European backward step \eqref{eq:european-step} can
be rewritten as
\begin{equation}
V_j^n = e^{-r\Delta t}p_u\,V_{j+1}^{n+1} + e^{-r\Delta t}p_m\,V_j^{n+1} + e^{-r\Delta t}p_d\,V_{j-1}^{n+1}.
\label{eq:trinomial-rewritten}
\end{equation}
Compare with the FTCS update for the transformed BS PDE
$V_\tau = (\sigma^2/2)V_{xx} + \nu V_x - rV$ (Block 1, with $\Delta\tau = \Delta t$):
\begin{equation}
V_j^{n+1} = (\alpha + \beta)V_{j+1}^n + (1 - 2\alpha - r\Delta t)V_j^n + (\alpha - \beta)V_{j-1}^n
\label{eq:ftcs}
\end{equation}
where $\alpha = (\sigma^2/2)\Delta t/\Delta x^2$ and
$\beta = \nu\Delta t/(2\Delta x)$.

\paragraph{Note.} Equation \eqref{eq:ftcs} marches forward in $\tau$
(and hence backward in $t$ since $\tau = T - t$), so the time index
on the LHS is $n+1$ while the RHS uses $n$. Equation
\eqref{eq:trinomial-rewritten} marches in calendar time, with the
LHS at $n$ and RHS at $n+1$. Both schemes go from later to earlier in
calendar time. The structural form is the same: a three-point linear
combination of the values one time level away.

\paragraph{Coefficient identification.} Substituting
$\Delta x = \sigma\sqrt{\lambda\Delta t}$ into the FTCS coefficients:
\begin{align*}
\alpha \;&=\; \frac{\sigma^2 \Delta t/2}{\lambda\,\sigma^2\,\Delta t} \;=\; \frac{1}{2\lambda},\\
\beta  \;&=\; \frac{\nu\,\Delta t}{2\,\sigma\sqrt{\lambda\,\Delta t}} \;=\; \frac{\nu\,\sqrt{\Delta t}}{2\,\sigma\,\sqrt{\lambda}}.
\end{align*}
Then the FTCS coefficients become
\begin{align*}
\alpha + \beta \;&=\; \frac{1}{2\lambda} + \frac{\nu\sqrt{\Delta t}}{2\sigma\sqrt{\lambda}} \;=\; p_u,\\
\alpha - \beta \;&=\; \frac{1}{2\lambda} - \frac{\nu\sqrt{\Delta t}}{2\sigma\sqrt{\lambda}} \;=\; p_d,\\
1 - 2\alpha - r\Delta t \;&=\; 1 - \frac{1}{\lambda} - r\Delta t \;=\; p_m - r\Delta t.
\end{align*}
The \emph{un-discounted} off-diagonal coefficients agree exactly:
$p_u = \alpha + \beta$ and $p_d = \alpha - \beta$ as identities.
The two schemes differ in \emph{how the discount factor is
distributed} across the recurrence:
\begin{itemize}[leftmargin=2em]
\item Trinomial multiplies the entire bracketed expression by
$e^{-r\Delta t}$, applying the discount to all three terms uniformly.
\item FTCS embeds the discount only in the diagonal coefficient
($1 - 2\alpha - r\Delta t = p_m - r\Delta t$), leaving the
off-diagonals undiscounted.
\end{itemize}
The middle coefficient comparison:
$e^{-r\Delta t}\,p_m \,-\, (p_m - r\Delta t) = r\Delta t / \lambda + O(\Delta t^2)$,
which is $O(\Delta t)$. The off-diagonal comparison
$e^{-r\Delta t}\,p_u - (\alpha + \beta) = -r\Delta t\,p_u + O(\Delta t^2)$
is also $O(\Delta t)$. The two schemes therefore agree at the leading
order $O(1)$ and differ at $O(\Delta t)$, which is the same order
as the local truncation error of both schemes. They are equivalent
in the sense that they share the same order of convergence and the
same limiting solution, but their numerical outputs at finite
$\Delta t$ differ by a small ($O(\Delta t)$) amount.

\paragraph{Conclusion.}
\begin{equation}
\boxed{\;\;\text{Trinomial KR with stretching } \lambda \;\sim\; \text{FTCS with } \alpha = \tfrac{1}{2\lambda} \;\;}
\label{eq:equivalence}
\end{equation}
in the asymptotic sense: both schemes have $O(\Delta t)$ local
truncation error, both converge to the same BS solution at the same
rate, and their numerical outputs differ by $O(\Delta t)$ (the same
order as the error itself). The two methods are not literally
identical, but they are structurally the same scheme written in two
languages.

\paragraph{The CFL constraint, in tree language.} Equation
\eqref{eq:equivalence} also explains why $\lambda \ge 1$ is a hard
constraint: $\lambda = 1$ corresponds to $\alpha = 1/2$, the FTCS
stability boundary. Going below ($\lambda < 1 \iff \alpha > 1/2$)
violates the CFL condition. In tree language this manifests as
$p_m < 0$: a probability that is no longer a probability. The same
phenomenon, two different symptoms.

\paragraph{Why $\lambda = 3$ is called ``third-order accurate''.}
With $\lambda = 3$, the probabilities are $(p_u, p_m, p_d) = (1/6, 2/3, 1/6) +
O(\sqrt{\Delta t})$. These are the weights of \emph{Simpson's rule}
for numerical integration. Simpson exactly integrates polynomials up
to degree~$3$. As a consequence, the third moment of the discrete
trinomial step matches the third moment of the GBM step (which is
$3\sigma^2\nu\Delta t^2 + O(\Delta t^3)$). Schemes that match higher
moments have smaller \emph{constants} in their error expansion, but
the global convergence rate remains $O(\Delta t)$ because the
non-smooth payoff (the kink at the strike) limits the achievable
order regardless. ``Third-order'' refers to moment matching, not
to global convergence.

% =====================================================================
\section{Convergence and comparison}
% =====================================================================

The trinomial converges to the BS price for European options, and
to the true American price for American options, at rate
$O(1/N)$ in the number of steps. This is the same rate as CRR
binomial. The kink-related limitation discussed in Block~3 (the
payoff is not $C^1$) caps both lattice methods at first order
regardless of the underlying scheme accuracy.

\begin{center}
\begin{tabular}{lll}
\hline
Method & Cost (American) & Convergence \\
\hline
CRR binomial         & $O(N^2)$ & $O(1/N)$ \\
Trinomial KR         & $O(N^2)$ & $O(1/N)$, smaller constant at $\lambda = 3$ \\
CN-PSOR              & $O(N^3 \log\epsilon^{-1})$ & $O(1/N)$ near the kink \\
\hline
\end{tabular}
\end{center}

The binomial and trinomial have the same asymptotic order but
trinomial typically gives slightly smaller errors at fixed $N$
(thanks to the third-moment matching). CN-PSOR has the highest order
but the highest cost. Block~6 will compare them numerically on a
common American put benchmark.

% =====================================================================
\section{Implementation outline}
% =====================================================================

A single module is enough. \texttt{quantlib.trinomial} (Python) and
\texttt{quant::pde::trinomial} (C++) provide:

\begin{itemize}[leftmargin=2em]
\item \texttt{trinomial\_european\_call(S, K, r, sigma, T, *, n\_steps, lambda\_param=3.0)}
\item \texttt{trinomial\_european\_put(S, K, r, sigma, T, *, n\_steps, lambda\_param=3.0)}
\item \texttt{trinomial\_american\_put(S, K, r, sigma, T, *, n\_steps, lambda\_param=3.0)}
\end{itemize}

Why no \texttt{trinomial\_american\_call}: for a stock with no
dividends, the American call price equals the European call price
(Merton 1973), so the function would duplicate the European version.

The internal structure is a single helper function \texttt{\_trinomial\_backward}
that performs the backward induction, parametrised by an option-type
flag (\texttt{european\_call}, \texttt{european\_put}, \texttt{american\_put})
that determines the terminal payoff and whether the early-exercise
$\max\{\cdot, \phi\}$ check is applied at each step.

% =====================================================================
\section{Validation strategy}
% =====================================================================

Five tests:

\begin{enumerate}[leftmargin=2em]
\item \textbf{Cross-validation (European)} against Black--Scholes
closed form on a parameter grid. Tolerance $\sim 10^{-3}$ at
$N = 2000$.

\item \textbf{Cross-validation (American)} against an embedded
CRR binomial reference (the same approach as Block~4's test). Tolerance
$\sim 5 \times 10^{-3}$.

\item \textbf{First-order convergence in $N$}: doubling $N$ should
halve the error, distinguishing the trinomial from any
hypothetical second-order scheme.

\item \textbf{Lambda sweep}: at fixed $N$, varying $\lambda \in
\{1.25, 2, 3, 5\}$ produces prices that all approach the BS limit
as $N$ grows. The $\lambda = 1$ degenerate case (binomial-equivalent)
agrees with the rest. This confirms that the choice of $\lambda$
within the valid range $\lambda \ge 1$ is a tuning preference,
not a correctness condition.

\item \textbf{Input validation}: $\lambda < 1$ raises (it would
violate $p_m \ge 0$); $n\_steps < 1$ raises; non-positive parameters
raise.
\end{enumerate}

The fourth test is the operational confirmation that the
$\lambda$-parameter is a tuning knob within a valid range, not a
correctness-determining parameter.

% =====================================================================
\section*{References}
% =====================================================================

\begin{thebibliography}{9}

\bibitem{KR1991}
B. Kamrad and P. Ritchken.
\emph{Multinomial approximating models for options with $k$ state variables.}
Management Science 37(12), 1640--1652, 1991. The original paper
introducing the $\lambda$-parametrised trinomial.

\bibitem{Hull2018}
J.~C. Hull.
\emph{Options, Futures, and Other Derivatives.}
10th ed., Pearson, 2018. Chapter 21: practitioner-oriented treatment
of binomial and trinomial trees.

\bibitem{Wilmott2006}
P. Wilmott.
\emph{Paul Wilmott on Quantitative Finance.}
2nd ed., Wiley, 2006. Chapter 20 on tree methods, including
implementation pseudocode.

\bibitem{TavellaRandall2000}
D. Tavella and C. Randall.
\emph{Pricing Financial Instruments: The Finite Difference Method.}
Wiley, 2000. Chapter 5 explicitly relates lattice methods to
finite-difference schemes; the FTCS--trinomial equivalence is
implicit in their treatment.

\bibitem{BrennanSchwartz1978}
M.~J. Brennan and E.~S. Schwartz.
\emph{Finite difference methods and jump processes arising in the
pricing of contingent claims: A synthesis.}
Journal of Financial and Quantitative Analysis 13(3), 461--474, 1978.
Historical reference for the equivalence between FD methods and
tree methods.

\bibitem{Merton1973}
R.~C. Merton.
\emph{Theory of rational option pricing.}
Bell Journal of Economics and Management Science 4(1), 141--183, 1973.
Includes the result that the American call equals the European call
for non-dividend-paying stocks.

\end{thebibliography}

\end{document}
